\chapter{Coding a Transformer from Scratch}
\label{chap:coding_scratch}

The best way to demystify Large Language Models is to build one. In this chapter, we will implement a small, character-level Generative Pre-trained Transformer (GPT) from scratch using PyTorch. This is heavily inspired by Andrej Karpathy's "nanoGPT."

\section{Prerequisites}

Ensure you have PyTorch installed:
\begin{lstlisting}[language=bash]
pip install torch
\end{lstlisting}

\section{Data Preparation \& Tokenization}

Before feeding text into a model, we must turn it into numbers.

\subsection{The Dataset}
For this example, we will use the "Tiny Shakespeare" dataset.

\begin{lstlisting}[language=Python]
# Download the dataset
!wget https://raw.githubusercontent.com/karpathy/char-rnn/master/data/tinyshakespeare/input.txt

with open('input.txt', 'r', encoding='utf-8') as f:
    text = f.read()

print("Length of dataset in characters: ", len(text))
# Output: ~1 million characters
\end{lstlisting}

\subsection{Character-Level Tokenizer}
We will build a simple character-level tokenizer. It maps every unique character to an integer.

\begin{lstlisting}[language=Python]
# unique characters
chars = sorted(list(set(text)))
vocab_size = len(chars)
print(''.join(chars))
print(vocab_size)

# Mapping from chars to ints and vice versa
stoi = { ch:i for i,ch in enumerate(chars) }
itos = { i:ch for i,ch in enumerate(chars) }
encode = lambda s: [stoi[c] for c in s] # encoder: take a string, output a list of integers
decode = lambda l: ''.join([itos[i] for i in l]) # decoder: take a list of integers, output a string

print(encode("hii there"))
# Output: [46, 47, 47, 1, 58, 46, 43, 56, 43]
print(decode(encode("hii there")))
# Output: "hii there"
\end{lstlisting}

\subsection{Byte Pair Encoding (BPE)}
Modern LLMs (GPT-4, Llama) use Byte Pair Encoding (BPE). BPE iteratively merges the most frequent pair of adjacent tokens.
\begin{enumerate}
    \item Start with individual characters: \texttt{['h', 'u', 'g', ' ', 'p', 'u', 'g']}
    \item Count frequency of pairs. "u" and "g" appear together often.
    \item Merge "u" + "g" -> "ug".
    \item New vocabulary: \texttt{['h', 'ug', ' ', 'p', 'ug']}
\end{enumerate}
This allows the model to handle common words as single tokens (efficient) while falling back to characters for rare words. For this tutorial, we stick to character-level for simplicity.

\subsection{Train/Val Split}
We split the data to ensure the model isn't just memorizing.

\begin{lstlisting}[language=Python]
import torch
data = torch.tensor(encode(text), dtype=torch.long)
n = int(0.9*len(data))
train_data = data[:n]
val_data = data[n:]

def get_batch(split):
    # generate a small batch of data of inputs x and targets y
    data = train_data if split == 'train' else val_data
    ix = torch.randint(len(data) - block_size, (batch_size,))
    x = torch.stack([data[i:i+block_size] for i in ix])
    y = torch.stack([data[i+1:i+block_size+1] for i in ix])
    x, y = x.to(device), y.to(device)
    return x, y
\end{lstlisting}

\section{The Architecture Components}

We will build the model bottom-up: starting with the smallest unit (Self-Attention Head) and assembling the full Transformer block.

\subsection{Imports and Hyperparameters}

First, let's set up our environment and define the model configuration.

\begin{lstlisting}[language=Python]
import torch
import torch.nn as nn
from torch.nn import functional as F

# Hyperparameters
batch_size = 32 # How many independent sequences to process in parallel?
block_size = 8  # What is the maximum context length for predictions?
max_iters = 3000
eval_interval = 300
learning_rate = 1e-3
device = 'cuda' if torch.cuda.is_available() else 'cpu'
eval_iters = 200
n_embd = 32
n_head = 4
n_layer = 4
dropout = 0.0

torch.manual_seed(1337)
\end{lstlisting}

\section{Self-Attention Head}

The heart of the Transformer. This module computes the attention scores between tokens.

\begin{lstlisting}[language=Python]
class Head(nn.Module):
    """ one head of self-attention """

    def __init__(self, head_size):
        super().__init__()
        self.key = nn.Linear(n_embd, head_size, bias=False)
        self.query = nn.Linear(n_embd, head_size, bias=False)
        self.value = nn.Linear(n_embd, head_size, bias=False)
        self.register_buffer('tril', torch.tril(torch.ones(block_size, block_size)))

        self.dropout = nn.Dropout(dropout)

    def forward(self, x):
        B,T,C = x.shape
        k = self.key(x)   # (B,T,C)
        q = self.query(x) # (B,T,C)

        # Compute attention scores ("affinities")
        wei = q @ k.transpose(-2,-1) * C**-0.5 # (B, T, C) @ (B, C, T) -> (B, T, T)
        wei = wei.masked_fill(self.tril[:T, :T] == 0, float('-inf')) # (B, T, T)
        wei = F.softmax(wei, dim=-1) # (B, T, T)
        wei = self.dropout(wei)

        # Perform the weighted aggregation of the values
        v = self.value(x) # (B,T,C)
        out = wei @ v # (B, T, T) @ (B, T, C) -> (B, T, C)
        return out
\end{lstlisting}

\section{Multi-Head Attention}

Multiple heads attending to different parts of the input in parallel.

\begin{lstlisting}[language=Python]
class MultiHeadAttention(nn.Module):
    """ multiple heads of self-attention in parallel """

    def __init__(self, num_heads, head_size):
        super().__init__()
        self.heads = nn.ModuleList([Head(head_size) for _ in range(num_heads)])
        self.proj = nn.Linear(n_embd, n_embd)
        self.dropout = nn.Dropout(dropout)

    def forward(self, x):
        out = torch.cat([h(x) for h in self.heads], dim=-1)
        out = self.dropout(self.proj(out))
        return out
\end{lstlisting}

\section{Feed-Forward Network}

A simple MLP applied to each position independently.

\begin{lstlisting}[language=Python]
class FeedFoward(nn.Module):
    """ a simple linear layer followed by a non-linearity """

    def __init__(self, n_embd):
        super().__init__()
        self.net = nn.Sequential(
            nn.Linear(n_embd, 4 * n_embd),
            nn.ReLU(),
            nn.Linear(4 * n_embd, n_embd),
            nn.Dropout(dropout),
        )

    def forward(self, x):
        return self.net(x)
\end{lstlisting}

\section{Transformer Block}

Combining attention and feed-forward layers with residual connections.

\begin{lstlisting}[language=Python]
class Block(nn.Module):
    """ Transformer block: communication followed by computation """

    def __init__(self, n_embd, n_head):
        # n_embd: embedding dimension, n_head: the number of heads we'd like
        super().__init__()
        head_size = n_embd // n_head
        self.sa = MultiHeadAttention(n_head, head_size)
        self.ffwd = FeedFoward(n_embd)
        self.ln1 = nn.LayerNorm(n_embd)
        self.ln2 = nn.LayerNorm(n_embd)

    def forward(self, x):
        x = x + self.sa(self.ln1(x))
        x = x + self.ffwd(self.ln2(x))
        return x
\end{lstlisting}

\section{The GPT Model}

Finally, the full model class.

\begin{lstlisting}[language=Python]
class GPTLanguageModel(nn.Module):

    def __init__(self):
        super().__init__()
        # each token directly reads off the logits for the next token from a lookup table
        self.token_embedding_table = nn.Embedding(vocab_size, n_embd)
        self.position_embedding_table = nn.Embedding(block_size, n_embd)
        self.blocks = nn.Sequential(*[Block(n_embd, n_head=n_head) for _ in range(n_layer)])
        self.ln_f = nn.LayerNorm(n_embd) # final layer norm
        self.lm_head = nn.Linear(n_embd, vocab_size)

    def forward(self, idx, targets=None):
        B, T = idx.shape

        # idx and targets are both (B,T) tensor of integers
        tok_emb = self.token_embedding_table(idx) # (B,T,C)
        pos_emb = self.position_embedding_table(torch.arange(T, device=device)) # (T,C)
        x = tok_emb + pos_emb # (B,T,C)
        x = self.blocks(x) # (B,T,C)
        x = self.ln_f(x) # (B,T,C)
        logits = self.lm_head(x) # (B,T,vocab_size)

        if targets is None:
            loss = None
        else:
            B, T, C = logits.shape
            logits = logits.view(B*T, C)
            targets = targets.view(B*T)
            loss = F.cross_entropy(logits, targets)

        return logits, loss

    def generate(self, idx, max_new_tokens):
        # idx is (B, T) array of indices in the current context
        for _ in range(max_new_tokens):
            # crop idx to the last block_size tokens
            idx_cond = idx[:, -block_size:]
            # get the predictions
            logits, loss = self(idx_cond)
            # focus only on the last time step
            logits = logits[:, -1, :] # becomes (B, C)
            # apply softmax to get probabilities
            probs = F.softmax(logits, dim=-1) # (B, C)
            # sample from the distribution
            idx_next = torch.multinomial(probs, num_samples=1) # (B, 1)
            # append sampled index to the running sequence
            idx = torch.cat((idx, idx_next), dim=1) # (B, T+1)
        return idx
\end{lstlisting}

\section{Training Loop}

The training loop is standard PyTorch boilerplate.

\begin{lstlisting}[language=Python]
model = GPTLanguageModel()
m = model.to(device)
optimizer = torch.optim.AdamW(model.parameters(), lr=learning_rate)

for iter in range(max_iters):

    # every once in a while evaluate the loss on train and val sets
    if iter % eval_interval == 0:
        losses = estimate_loss()
        print(f"step {iter}: train loss {losses['train']:.4f}, val loss {losses['val']:.4f}")

    # sample a batch of data
    xb, yb = get_batch('train')

    # evaluate the loss
    logits, loss = model(xb, yb)
    optimizer.zero_grad(set_to_none=True)
    loss.backward()
    optimizer.step()

# generate from the model
context = torch.zeros((1, 1), dtype=torch.long, device=device)
print(decode(m.generate(context, max_new_tokens=500)[0].tolist()))
\end{lstlisting}

\section{Conclusion}

This implementation is simplified but structurally identical to GPT-2. The magic of Deep Learning is that scaling this simple architecture---more layers, more heads, bigger embedding dimensions, and much more data---results in the emergent capabilities we see in modern LLMs.

You have now moved from a consumer of AI to a builder of AI. The journey continues with experimentation: try changing the activation functions, add rotational embeddings, or train it on different datasets (Shakespeare, Python code, etc.).
